\section{问题重述}

\subsection{问题背景}

随着全球贸易的快速发展，集装箱作为国际货物运输的核心载体，被广泛应用于海运、陆运及多式联运环节。长期的装卸、堆叠与运输过程中，集装箱经常受到机械冲击、腐蚀、变形等外力作用，导致箱体外表面产生裂纹、凹陷、穿孔、锈蚀等不同类型的损伤。这些破损不仅影响集装箱的结构强度与密封性能，造成货物泄漏、污染或丢失，还可能引发运输安全事故并造成经济损失。因此，实现对集装箱破损的高效、准确、标准化检测，已成为港口运输管理与国际物流安全中的关键问题。

传统的集装箱破损检测主要依赖人工巡检或简单图像比对，存在检测效率低、主观性强、环境适应性差等问题，难以满足港口自动化与大规模作业的需求。近年来，随着计算机视觉与深度学习技术的迅速发展，基于图像识别的自动检测方法成为研究热点。通过固定摄像头或龙门吊摄像头自动采集集装箱图像，并引入卷积神经网络（CNN\cite{he2016resnet}）与目标检测模型（如 YOLO 系列\cite{redmon2016,redmon2018,bochkovskiy2020,jocher2023}、Faster R-CNN\cite{ren2015} 等），可在复杂背景下自动识别多类型破损，实现集装箱状态的智能化评估与实时监测，为智慧港口建设和安全运输提供技术支撑。

\subsection{问题提出}

本题要求利用提供的集装箱外表面图像数据集，构建智能识别模型，完成以下三项任务：

\begin{enumerate}[leftmargin=2.2em,itemsep=2pt]
  \item {\heiti 问题一（分类任务）：}结合给出的图像文件，提取图像特征，建立模型判断该集装箱图像是否存在任何形式的残损；
  \item {\heiti 问题二（检测与分割任务）：}如果存在残损，需要定位残损所在位置，并识别残损的具体类别；
  \item {\heiti 问题三（模型评估）：}对上述问题一、问题二中建立的模型，从不同维度进行评估。
\end{enumerate}

此外，需要利用所建立的模型对测试集中的图片进行检测，并将检测结果按照 \texttt{image\_id,\allowbreak class\_id,\allowbreak x\_center,\allowbreak y\_center,\allowbreak width,\allowbreak height} 的格式写入 \texttt{test\_result.csv} 中，其中坐标均为归一化值。

赛题同时提示了数据与建模中的主要困难：图像背景复杂（港口机械、天空、地面），存在光照变化、反光、阴影、雨水、污渍等干扰；残损大小不一，大的如大面积锈蚀，小的如细微裂缝，对模型的多尺度特征提取能力要求高；无残损图片可能远多于有残损图片；某些残损类别（如裂缝）的样本数量可能远少于其他类别（如凹陷）；某些残损类别外观相似（如深凹陷与破损），难以区分。上述提示构成了本文建模与实验设计的重要出发点。
